\input{myquizpreamble}
\input{yliow}
\input{ciss362}
\textwidth=6in

\renewcommand\TITLE{Assignment a10}


\renewcommand\AUTHOR{John Doe}

\begin{document}
\topmatter

\textsc{Objectives}
\begin{itemize}
\li Design context-free grammars.
\li Design PDA.
\li Prove a language is not context-free.
\end{itemize}

For questions asking you show an operation on context-free languages
language is closed, you need not prove your construction is correct.
(But you are welcome to do so. Most of such proofs, if not immediate,
is by induction.)
If the question does not say so, you can choose to give a PDA or a
context-free grammar.
For a PDA, draw the PDA state diagram and give the formal definition.

Study the solutions to the following questions in the textbook:
2.18, 2.30(b) and (c)

Here are the questions you should work.
Some questions have solutions (either from
the book or I have written up the solution).
\begin{myenum}
\li Sipser 2.11. Q1. 
\li Sipser 2.12. Q2. 
\li Sipser 2.14. Q3. 
\li Sipser 2.15. Solution to 2.15 is provided.
\li Sipser 2.16. Q4.
\li Sipser 2.17. Solution to 2.17 is provided.
\li Sipser 2.18. I talked about this in class. Solution to 2.18 is also provided in the textbook -- study it carefully.
\li Sipser 2.19. Q5. Solutions to 2.6(a) and 2.6(c) are provided in the textbook. Study it carefully.
\li Sipser 2.20. Q6.
\li Sipser 2.30. Q7. 
\li Sipser 2.31. Q8.
\li Sipser 2.32. Q9.
\li Q10.  
\end{myenum}

%------------------------------------------------------------------------------
\newpage
\textsc{How to draw a state diagram}

Here's an example showing you how to draw the elements of a state diagram.
Also, look at the solution to 1.3 below.

\begin{center}
  \begin{tikzpicture}[shorten >=1pt,>=triangle 60,double distance=2pt,node distance=2cm,auto,initial text=]
    
    \node[state,initial]   (q0) at (0, 0) {$q_0$};
    \node[state]           (q1) at (4, 0) {$q_1$};
    \node[state,accepting] (q2) at (8, 0) {$q_2$};
    \node[state]           (q3) at (0,-4) {$q_3$};
    \node[state]           (q4) at (4,-4) {$q_4$};

    \node[state]           (q5) at (12, 0) {$q_5$};
    \node[state]           (q6) at (12,-4) {$q_6$};

    \path[->]
    (q0) edge                node {$1$} (q1)    
    (q1) edge [bend left=15] node {$0$} (q2)
    (q0) edge [left]  node {$\alpha$} (q3)
    (q1) edge [right] node {$\beta$}  (q4)
    (q1) edge [loop above]   node {$1$} (q1)
    (q2) edge [loop below]   node {$0$} ()
    (q2) edge [bend left=15] node {$1$} (q1)
    (q3) edge [loop right]   node {$0,1$} ()
    (q4) edge [loop left]    node {$2,3$} ()
    (q5) edge [right, bend left=30]    node {$\gamma$} (q6)
    (q5) edge [right, bend left=60]    node {$\ep$}    (q6)
    (q5) edge [left, bend right=30]    node {$\delta$} (q6)
    ;
  \end{tikzpicture}
\end{center}

For more information on drawing state diagrams go to my tutorials and
look for \verb!latex-automata.pdf!:
\begin{center}
{\small \url{https://drive.google.com/file/d/1AeE-P0WNvQlitzPDxQpGE8bMR9Yc9gMW}}
\end{center}
Let me know if you have any questions about drawing state diagram.


%------------------------------------------------------------------------------
\newpage
\textsc{How to write a context-free grammar}

Here's a context-free grammar, $G$, for our $\{ a^n b^n \mid n \geq 0 \}$:
\[
G:
\begin{cases}
S \rightarrow \ep \\
S \rightarrow aSb
\end{cases}
\]
Usually we combine the productions for the same variable like this:
\[
G:
S \rightarrow \ep \mid aSb
\]
Here's an example of a longer grammar:
\[
G_1:
\begin{cases}
S \rightarrow \ep \mid TU \\
T \rightarrow a \mid   aT \\
U \rightarrow \ep \mid bUcc
\end{cases}
\]
And remember to list the starting variable first.

%------------------------------------------------------------------------------
\newpage

\textsc{How to draw a CYK table}

\begin{python}
from latextool_basic import *
p = Plot()
m = [['$\{B,C\}$',  '$\{A\}$',     '$\{A\}$', '$\{A\}$',  '$\{A\}$',  '$\{B,C\}$'],
     ['$\emptyset$','$\{C\}$',     '$\{C\}$', '$\{C\}$',  '$\{S\}$',  ''         ],
     ['$\{A\}$',    '$\{S\}$',     '$\{S\}$', '$\{B\}$',  '',         ''         ],
     ['$\{C\}$',    '$\emptyset$', '$\{S\}$', '',         '',         ''         ],
     ['$\{S\}$',    '$\{B\}$',     '',        '',         '',         ''         ],
     ['$\{B,S\}$',  '',            '',        '',         '',         ''         ],
     ]
cyk(p, m, w='baaaab',fontsize='small')
print(p)
\end{python}

The following shows you how to control the size of the cells:
\begin{python}
from latextool_basic import *
p = Plot()
m = [[r'$\{B_b,C_b\}$',         r'$\{A_a\}$',           r'$\{A_a\}$',           r'$\{A_a\}$',           r'$\{A_a\}$',           r'$\{B_b,C_b\}$'],
     [r'$\emptyset$',           r'$\{C_{(1,2),(1,3)}\}$', r'$\{C_{(1,3),(1,4)}\}$', r'$\{C_{(1,4),(1,5)}\}$', r'$\{S_{(1,5),(1,6)}\}$', r''],
     [r'$\{A_{(1,1),(2,2)}\}$',   r'$\{S_{(1,2),(2,3)}\}$', r'$\{S_{(1,3),(2,4)}\}$', r'$\{B_{(2,4),(1,6)}\}$', r'',                    r''],
     [r'$\{C_{(3,1),(1,4)}\}$',   r'$\emptyset$',         r'$\{S_{(1,3),(3,4)}\}$', r'',                   r'',                    r''],
     [r'$\{S_{(3,1),(2,4)}\}$',   r'$\{B_{(2,2),(3,4)}\}$', r'',                    r'',                   r'',                    r''],
     [r'$\{B_{(1,1),(2,5)}$ $S_{(3,1),(3,4)}\}$', r'',      r'',                    r'',                   r'',                    r''],
     ]
def getrect():
    width = 2.6  # width of cell
    height = 1.4 # height of cell
    def rect(x):
        return Rect(x0=0, y0=0, x1=width, y1=height,
                    innersep=0.2,
                    s='%s' % x, align='t')
    return rect
    
cyk(p, m, w='baaaab', fontsize='small', rect=getrect())
print(p)
\end{python}

%------------------------------------------------------------------------------
\newpage
Q1. Sipser 2.11. 

\textsc{Solution}.

\input{q01.tex}

%------------------------------------------------------------------------------
\newpage
Q2. Sipser 2.12.

\textsc{Solution}.

\input{q02.tex}

%------------------------------------------------------------------------------
\newpage
Q3. Sipser 2.14.

\textsc{Solution}.

\input{q03.tex}

%------------------------------------------------------------------------------
\newpage
Sipser 2.15.

Give a counter-example to show that the following fails to prove
that the class of context-free languages is closed under Kleene
star.

\lq\lq
Let $L$ be a CFL that is generated by the CFG
$G = (\Sigma, V, S, P)$ where $V$ is the set of variables,
$S \in V$ is the starting symbol, and $P$ is the set of production rules.
Add the new rule $S \rightarrow SS$ and call the resulting
grammar $G'$.
$G'$ generates $L^*$."

\textsc{Solution}.

Consider the context-free grammar $G: S \rightarrow a$.
Then $L(G')$ does not contain $\ep$.
(In fact $L(G') = \{a^n \mid n \geq 2\}$.)
\qed

%------------------------------------------------------------------------------
\newpage
Q4. Sipser 2.16

Show that the class of context-free languages is closed under
union, concatenation, and Kleene star.

(Note that I have actually already given you the proofs in class.
You just have to write it up in proper mathematical prose.)

\textsc{Solution}.

\input{q04.tex}

%------------------------------------------------------------------------------
\newpage
Sipser 2.17.
Use the results of Exercise 2.16 to give another proof that
every regular language is context free, by showing how to convert
a regular expression directly to an equivalent context-free
grammar.

\textsc{Solution}.

We will prove by induction on the left of the regular expression $r$
such that $L(r)$ is context-free.
Let $P(n)$ be the statement that if $r$ is a regular expression
of length $n$, then $L(r)$ is context-free.

\textsc{Base case}.
Let $r$ be a regular expression of length $1$.
Then either $r = c \in \Sigma$ or $r = \emptyset$.
If $c \in \Sigma$, then $L(c) = \{c\}$ is context-free since
the grammar $S \rightarrow c$ generated $L(c)$.
Also, the grammar $S \rightarrow S$ generates $L(\emptyset)$.
Therefore $P(1)$ holds.

\textsc{Inductive case}.
Assume $P(k)$ holds for $1 \leq k \leq n$.
We will show that $P(n + 1)$ holds where $n + 1 \geq 2$.
Let $r$ be a regular expression of length $n + 1$.
Then $r$ is $(r_0), r_0^*, r_0 \cup r_1, r_0 \cdot r_1$.
In all cases $r_0$ or $r_0, r_1$ are regular expressions of length $\leq n$.
Hence $L(r_0)$ or $L(r_0), L(r_1)$ are context-free.
Since context-free languages are closed under concatenation, union, and Kleene star,
$L(r)$ is also context-free.
(The case of $r = (r_0)$ is immediate and does not require closure of concatention, union, Kleene star.)
Hence $P(n + 1)$ holds.

Therefore by strong mathematical induction $P(n)$ holds for all $n \geq 1$, i.e.,
all regular languages are context-free.
\qed

%------------------------------------------------------------------------------
\newpage
Q5. Sipser 2.19.

\textsc{Solution}.

\input{q05.tex}

%------------------------------------------------------------------------------
\newpage
Q6. Sipser 2.20.

\textsc{Solution}.

\input{q06.tex}

%------------------------------------------------------------------------------
\newpage
Q7. Sipser 2.21.
(You need not provide the proof.)

\textsc{Solution}.

\input{q07.tex}

%------------------------------------------------------------------------------
\newpage
Q8. Sipser 2.22.

\textsc{Solution}.

\input{q08.tex}

%------------------------------------------------------------------------------
\newpage
Q9. Sipser 2.32.

\textsc{Solution}.

\input{q08.tex}

%------------------------------------------------------------------------------

Q10.
Consider the grammar $G$ given by
  \begin{align*}
    S &\rightarrow AB \mid BA \\
    A &\rightarrow AC \mid a \\
    B &\rightarrow SA \mid SB \mid b \\
    C &\rightarrow SS \mid BA \mid a \\
  \end{align*}
  Using the grammar $G$ above, check if the following words are generated
  by $G$ using the CKY algorithm. Provide the completed CYK tables.
  For words generated by $G$, write down all their
  possible derivations. 
  \begin{enumerate}
  \item[(a)] $aa$
  \item[(b)] $bab$
  \item[(c)] $abab$
  \item[(d)] $aabaab$
  \end{enumerate}

\textsc{Solution}.

\input{q10.tex}



\end{document}
